\documentclass{article}
\usepackage{hyperref}

\title{nanoSoC Test Board Datasheet}
\author{\href{http://www.soclabs.org}{SoC Labs}}
\begin{document}
\maketitle

\tableofcontents
\clearpage
\section{Introduction}

\section{Summary}

\section{FPGA Connections}

\begin{table}[h!]
    \centering
    \begin{tabular}{ c c c c c}
    FPGA Pin Name & FPGA Pin Number & FPGA Bank &  Nanosoc Pin Name & Nanosoc Pin Number \\
    \hline
    PL16A   & 25    & 3    &  &   \\
    PL16B   & 26    & 3    &  &   \\
    PL17A   & 27    & 3    &  &   \\
    PL17B   & 28    & 3    &  &   \\
    PL18D   & 31    & 3    & SE     &   \\
    PL19A   & 32    & 3    & TEST   &   \\
    PL19B   & 33    & 3    & P0[0]  &   \\
    PL20A   & 34    & 3    & P0[1]  &   \\
    PL20B   & 35    & 3    & P0[2]  &   \\
    PB3A    & 38    & 2    & P0[3]  &   \\
    PB3B    & 39    & 2    & P0[4]  &   \\
    PB4A    & 40    & 2    & P0[5]  &   \\
    PB4B    & 41    & 2    & P0[6]  &   \\
    PB7A    & 42    & 2    & P0[7]  &   \\
    PB7B    & 43    & 2    & P0[8]  &   \\
    PB9A    & 44    & 2    & P0[9]  &   \\
    PB9B    & 45    & 2    & P0[10] &   \\
    PB10A   & 47    & 2    & P0[11] &   \\
    PB10B   & 48    & 2    & P0[12] &   \\
    PB13A   & 49    & 2    & P0[13] &   \\
    PB13B   & 50    & 2    & P0[14] &   \\
    PB15A   & 52    & 2    & P0[15] &   \\
    PB15B   & 54    & 2    & P1[0]  &   \\
    PB20A   & 55    & 2    & P1[1]  &   \\
    PB20B   & 56    & 2    & P1[2]  &   \\
    PB21A   & 57    & 2    & P1[3]  &   \\
    PB21B   & 58    & 2    & P1[4]  &   \\
    PB23A   & 59    & 2    & P1[5]  &   \\
    PB23B   & 60    & 2    & P1[6]  &   \\
    PB24A   & 61    & 2    & P1[7]  &   \\
    PB24B   & 62    & 2    & P1[8]  &   \\
    PB24D   & 63    & 2    & P1[9]  &   \\
    PB27A   & 65    & 2    & P1[10] &   \\
    PB27B   & 67    & 2    & P1[11] &   \\
    PB29A   & 68    & 2    & P1[12] &   \\
    PB29B   & 69    & 2    & P1[13] &   \\
    PB30A   & 70    & 2    & P1[14] &   \\
    PB30B   & 71    & 2    & P1[15] &   
    \end{tabular}
    \caption{FPGA Connections to nanoSoC}
    \label{table:fpga_nanosoc}
\end{table}


\begin{table}[h!]
    \centering
    \begin{tabular}{ c c c c c}
    FPGA Pin Name & FPGA Pin Number & FPGA Bank & Header Pin Name & Header Pin Number \\
    \hline
    PT28D   & 109   & 0    & SA3 &   \\
    PT28C   & 110   & 0    & SA2 &   \\
    PT27B   & 111   & 0    & SA1 &   \\
    PT27A   & 112   & 0    & SA0 &   \\
    PT25B   & 113   & 0    &    &   \\
    PT25A   & 114   & 0    &    &   \\
    PT24B   & 115   & 0    &   &   \\
    PT24A   & 117   & 0    &   &   \\
    PT23D   & 119   & 0    &   &   \\
    PT23C   & 120   & 0    &   &   \\
    PT21B   & 121   & 0    &   &   \\
    PT21A   & 122   & 0    &   &   \\
    PT20D   & 125   & 0    & SDA  &   \\
    PT20C   & 126   & 0    & SCL  &   \\
    PT18B   & 127   & 0    &   &   \\
    PT18A   & 128   & 0    &   &   \\
    PT15D   & 130   & 0    &  &   \\
    PT15C   & 131   & 0    &  &   \\
    PT14B   & 132   & 0    &  &   \\
    PT14A   & 133   & 0    &  &   \\
    PT13D   & 136   & 0    &  &   \\
    PT13C   & 137   & 0    &  &   \\
    PT11B   & 138   & 0    &   &   \\
    PT11A   & 139   & 0    &   &   \\
    PT10B   & 140   & 0    &   &   \\
    PT10A   & 141   & 0    &   &   \\
    PT9B    & 142   & 0    &   &   \\
    PT9A    & 143   & 0    &   &   \\
      &    & 2    &   &   \\
      &    & 2    &   &   \\
      &    & 2    &   &   \\
      &    & 2    &   &   \\
      &    & 2    &  &   \\
      &    & 2    &  &   \\
      &    & 2    &  &   \\
      &    & 2    &  &   \\
      &    & 2    &  &   \\
      &    & 2    &  &   
    \end{tabular}
    \caption{FPGA to Pi Header Connections}
    \label{table:fpga_pi}
\end{table}


\begin{table}[h!]
    \centering
    \begin{tabular}{ c c c c c}
    FPGA Pin Name & FPGA Pin Number & FPGA Bank & PMOD Pin \\
    \hline
    PL3A   & 1  & 5   & 8   \\
    PL3B   & 2  & 5   & 4   \\
    PL4A   & 3  & 5   & 7   \\
    PL4B   & 4  & 5   & 3   \\
    PL7A   & 9  & 5   & 6   \\
    PL7B   & 10 & 5   & 2   \\
    PL8A   & 11 & 5   & 5   \\
    PL8B   & 12 & 5   & 1   \\
    \end{tabular}
    \caption{FPGA to PMODA Connections}
    \label{table:fpga_pmoda}
\end{table}


\begin{table}[h!]
    \centering
    \begin{tabular}{ c c c c c}
    FPGA Pin Name & FPGA Pin Number & FPGA Bank & PMOD Pin \\
    \hline
    PL9A   & 13 & 4   & 8   \\
    PL9B   & 14 & 4   & 4   \\
    PL10A  & 15 & 4   & 7   \\
    PL10B  & 17 & 4   & 3   \\
    PL13A  & 21 & 4   & 6   \\
    PL13B  & 22 & 4   & 2   \\
    PL14A  & 23 & 4   & 5   \\
    PL14B  & 24 & 4   & 1   \\
    \end{tabular}
    \caption{FPGA to PMODB Connections}
    \label{table:fpga_pmodb}
\end{table}



\section{Communications interface}

\subsection{I2C}

\begin{center}
    \begin{tabular}{ c c c c}
    Device & Identifier & Adress & Purpose \\
    \hline
    CDCM6208 & U1 & 10101(AD1 AD0) & Clock Controller \\
    INA228 & U2 & 1000000 &  VDDACC Monitor \\
    INA228 & U3 & 1000001 & VDD CORE Monitor \\
    INA228 & U4 & 1000100 & VDDIO Monitor \\
    LP87563-Q1 & & 1100001 & Power Management IC 
    \end{tabular}
\end{center}

\section{References}
Bugblat PIF uses raspberry pi to program the Mach XO2 Chip: \href{https://www.bugblat.com/products/pif/index.html}{PIF link}

\end{document}